\begin{abstract}
Сошиал медиа мониторингийн систем нь сошиал орчин дахь хэрэглэгчдийн нийтлэл, пост, жиргээ, сэтгэгдэл зэрэг мэдээллийг бодит цагийн горимд нэг платформд төвлөрүүлэн цуглуулж, боловсруулан, ойлгомжтой хэлбэрээр харуулдаг цогц шийдэл юм. Ийм төрлийн систем нь маркетинг, олон нийттэй харилцах үйл ажиллагаа, аюулгүй байдал, судалгаа шинжилгээ зэрэг олон салбарт өргөн хэрэглэгдэж байна. Тус систем нь байгууллага болон хувь хүмүүст сошиал орчин дахь мэдээллийн урсгалыг үр дүнтэй удирдах, задлан шинжлэх, хянах боломжийг олгохоос гадна үйл явдлыг цаг алдалгүй илрүүлэх, нөлөөлөгчдийг тодорхойлох, чиг хандлагыг таних зэрэг чухал үүргийг гүйцэтгэдэг.

Sumbee.mn нь сошиал мониторингийн чиглэлээр монголд үйл ажиллагаа явуулж буй анхдагч компаниудын нэг бөгөөд одоогоор түлхүүр үг болон хиймэл оюунд (LLM) суурилсан тоон болон сэтгэгдлийн (sentiment) шинжилгээ хийх боломжтой системийг хөгжүүлэн ажиллаж байна. Энэхүү дипломын ажлын хүрээнд Sumbee.mn компанийн сошиал мониторингийн системийн архитектур, ашиглаж буй технологи, хэрэглээний жишээ кейсүүдийг судалж, системийн үр ашиг, гүйцэтгэлийг нэмэгдүүлэх боломжит шийдлүүдийг тодорхойлно.

Одоогийн байдлаар Sumbee.mn систем дэх их өгөгдлийн шинжилгээ нь тухайн сэдвийн хүрээнд ямар агуулга давамгайлан яригдаж байгаа, түүнээс ямар нарийвчилсан дүгнэлт гаргаж болохыг гүнзгий түвшинд тодорхойлох чадвар харьцангуй сул байгаа нь анхаарал хандуулах шаардлагатай асуудал болоод байна.

Иймээс энэхүү дипломын ажлын гол зорилго нь Sumbee.mn компанийн сошиал мониторингийн системийн өгөгдлийн шинжилгээний үр дүнг сайжруулах, хэрэглэгчдэд илүү нарийвчилсан, оновчтой дүгнэлт гаргах боломж бүрдүүлэх зорилгоор Topic Modeling аргачлалуудыг судалж, тохиромжтой шийдэл боловсруулахад оршино.

\section{Судалгааны үндэс ба асуудлын тодорхойлолт}
Сүүлийн жилүүдэд сошиал медиа нь мэдээлэл түгээх, олон нийтийн санал бодлыг илэрхийлэх, брэндийн имидж бүрдүүлэх, хэрэглэгчийн зан төлөвийг тодорхойлох чухал суваг болж байна. Үүнтэй холбоотойгоор сошиал орчноос их хэмжээний өгөгдөл тасралтгүй үүсч, байгууллагууд энэхүү их өгөгдлийг үр дүнтэй ашиглах шаардлага улам бүр нэмэгдэж байна.

Сошиал медиа мониторингийн системүүд нь энэхүү өгөгдлийг бодит цаг хугацаанд цуглуулж, ангилан боловсруулж, тоон болон сэтгэгдлийн (sentiment) шинжилгээ хийх боломжийг олгодог. Гэвч зөвхөн түлхүүр үг болон энгийн sentiment анализ дээр тулгуурласан шинжилгээ нь тухайн сэдвийн хүрээнд яг ямар дэд сэдвүүд, ямар гол агуулгууд давамгайлж байгааг гүнзгий түвшинд илрүүлэхэд хангалтгүй байдаг.

Sumbee.mn компанийн хувьд сошиал мониторингийн системийг хөгжүүлэн ашиглаж байгаа боловч их өгөгдлийн шинжилгээний үр дүн нь тухайн сэдвийн хүрээнд яригдаж буй агуулгыг нарийвчилсан түвшинд ангилан тодорхойлох, оновчтой дүгнэлт гаргах тал дээр харьцангуй сул байгаа нь тулгамдсан асуудал болж байна. Өөрөөр хэлбэл, одоогийн систем нь мэдээллийг тоон хэмжээнд болон эерэг, сөрөг, саармаг байдлаар ангилж чаддаг боловч тухайн өгөгдлийн гүн агуулгыг илрүүлж, бодитой шийдвэр гаргалтад ашиглахуйц бүтэцчилсэн мэдлэг болгон хувиргах боломж хязгаарлагдмал байна.

Иймд сошиал орчноос цугларсан их хэмжээний текстэн өгөгдлөөс сэдэв, дэд сэдвүүдийг автоматаар илрүүлэх Topic Modeling аргачлалыг судалж, практикт нэвтрүүлэх шаардлага үүсч байна. Энэ нь системийн аналитик чадамжийг нэмэгдүүлж, хэрэглэгчдэд илүү бодит, нарийвчилсан, шийдвэр гаргалтад дэмжлэг үзүүлэхүйц дүн шинжилгээ өгөх боломжийг бүрдүүлнэ.
\section{Судалгааны зорилго ба зорилт}
Энэхүү судалгааны ажлын зорилго нь Sumbee.mn компанийн сошиал мониторингийн системийн өгөгдлийн шинжилгээний үр дүнг сайжруулах зорилгоор Topic Modeling аргачлалуудыг судалж, тохиромжтой шийдэл боловсруулан системд нэвтрүүлэх боломжийг тодорхойлоход оршино. Үүний хүрээнд дараах зорилтуудыг дэвшүүлж байна. Үүнд:
\begin{enumerate}
    \item Сошиал медиа мониторингийн системийн онолын үндэс, архитектур, өгөгдлийн боловсруулалтын аргачлалуудыг судлах.
    \item Topic Modeling-ийн онолын суурь ойлголт, түгээмэл алгоритмууд (жишээлбэл, LDA, NMF, embedding-д суурилсан аргачлалууд гэх мэт)-ыг судалж, харьцуулсан дүн шинжилгээ хийх.
    \item Sumbee.mn системийн одоогийн өгөгдлийн шинжилгээний үйл явц, сул болон давуу талуудыг тодорхойлох.
    \item Сошиал орчноос цугларсан текстэн өгөгдөл дээр туршилт хийж, тохиромжтой Topic Modeling аргачлалыг сонгон хэрэгжүүлэх.
    \item Хэрэгжүүлсэн шийдлийн үр дүнг үнэлж, системийн аналитик чанар, хэрэглээний үр ашгийг сайжруулах боломжийг тодорхойлох.
\end{enumerate}
Эдгээр зорилтуудыг хэрэгжүүлснээр сошиал мониторингийн системийн өгөгдлийн шинжилгээний гүнзгийрэл нэмэгдэж, хэрэглэгчдэд илүү бодитой, бүтэцчилсэн, шийдвэр гаргалтад чиглэсэн мэдээлэл хүргэх нөхцөл бүрдэх юм.
\section{Судалгааны арга зүй, хүрээ ба хязгаарлалт}
\textbf{Судалгааны арга зүй}

Энэхүү судалгаанд чанарын (qualitative) болон тоон (quantitative) шинжилгээний аргуудыг хослуулан ашиглав. Судалгааны арга зүй нь онолын судалгаа, системийн архитектурын шинжилгээ, өгөгдөлд суурилсан туршилт, кейс судалгаанд тулгуурласан болно.

Онолын судалгааны хүрээнд сошиал медиа мониторингийн системийн бүтэц, их өгөгдлийн боловсруулалт, текстэн өгөгдлийн шинжилгээ, Topic Modeling аргачлалуудын онолын үндэс, түгээмэл алгоритмууд (LDA, NMF, embedding-д суурилсан аргачлалууд зэрэг)-ыг судалж, харьцуулсан дүн шинжилгээ хийв. Мөн орчин үеийн NLP болон LLM-д суурилсан текст боловсруулалтын чиг хандлагыг авч үзсэн.

Системийн архитектурын шинжилгээний хүрээнд Sumbee.mn сошиал мониторингийн системийн өгөгдөл цуглуулах, хадгалах, боловсруулах, шинжлэх одоогийн үйл явцыг задлан шинжилж, өгөгдлийн урсгал болон аналитик модулийн уялдаа холбоог тодорхойлсон. Одоогийн түлхүүр үг болон sentiment анализд суурилсан шийдлийн давуу болон сул талуудыг илрүүлсэн.

Өгөгдөлд суурилсан туршилтын хүрээнд сошиал орчноос цугларсан бодит текстэн өгөгдөл дээр Topic Modeling алгоритмуудыг хэрэгжүүлж, гарсан үр дүнг чанарын болон тоон үзүүлэлтээр үнэлэв. Сэдвийн уялдаа, ойлгомжтой байдал, практик хэрэглээний ач холбогдлыг харьцуулан дүгнэсэн.

Кейс судалгааны хүрээнд тодорхой нэг сэдэв, брэнд, эсвэл нийгмийн асуудлын хүрээнд цугларсан өгөгдөл дээр туршилт хийж, Topic Modeling нэвтрүүлснээр гарсан үр дүн нь бодит шийдвэр гаргалтад хэрхэн дэмжлэг үзүүлж байгааг үнэлсэн.
\textbf{Судалгааны хүрээ}

Судалгааны ажлын хүрээ нь Sumbee.mn компанийн сошиал мониторингийн системд төвлөрнө. Үүнд:

Сошиал орчноос цугларсан текстэн өгөгдлийн боловсруулалт

Түлхүүр үг болон sentiment анализд суурилсан одоогийн шинжилгээ

Topic Modeling аргачлалын онол, алгоритм, хэрэгжилт

Системийн аналитик модулийн сайжруулалт

Судалгаанд Монгол хэл дээрх сошиал медиа текстэн өгөгдлийг ашигласан бөгөөд тодорхой сэдэв, кампанит ажил, эсвэл байгууллагатай холбоотой бодит кейс дээр туршилт хийсэн. Судалгааны гол анхаарал нь өгөгдлийн гүн агуулгыг илрүүлэх, сэдэвчилсэн бүтэц гаргах, аналитик чанарыг нэмэгдүүлэхэд чиглэв.

\textbf{Судалгааны хязгаарлалт}

Энэхүү судалгаа нь дараах хязгаарлалтуудтай:
\begin{itemize}
    \item Судалгаа нь Sumbee.mn системийн хүрээнд хийгдсэн бөгөөд бусад сошиал мониторингийн системүүдтэй өргөн хүрээний туршилт, гүйцэтгэлийн харьцуулалт хийгдээгүй.
    \item Туршилтад ашигласан өгөгдөл нь тодорхой хугацаа, тодорхой сэдвийн хүрээнд цугларсан тул бүх төрлийн сошиал өгөгдлийг бүрэн төлөөлөхгүй байж болзошгүй.
    \item Topic Modeling алгоритмуудын гүйцэтгэлийн үнэлгээ нь тодорхой шалгуур үзүүлэлтүүдэд (сэдвийн уялдаа, ойлгомжтой байдал, практик ач холбогдол гэх мэт) тулгуурласан бөгөөд хэрэглэгчийн өргөн хүрээний бодит туршилтаар баталгаажуулалт хязгаарлагдмал хийгдсэн.
    \item Боловсруулсан шийдэл нь туршилтын (prototype) түвшинд хэрэгжсэн бөгөөд үйлдвэрлэлийн орчинд бүрэн хэмжээний нэвтрүүлэлт, урт хугацааны гүйцэтгэлийн үнэлгээ хийгдээгүй.
\end{itemize}
\end{abstract}
